% 张三慧第4版第二卷第28章《原子中的电子》章末习题28.1--28.34独立解析。
% 主教材PDF423--426，辅导S PDF373--383；版权/学术和原图重绘尚待审查。
% 仅候选：exam_ready=false。
\section*{第28章\quad 原子中的电子：独立解析}
\paragraph{习题 28.1}
Balmer谱系的电子终态 \(n_f=2\)，起态
\(n_i=3,4,\ldots\)，Rydberg公式
\(1/\lambda=R_H(1/4-1/n_i^2)\)。
最长波长为 \(n_i=3\to2\) 的 \(H_\alpha\)，
最短为系列极限 \(n_i\to\infty\)：
\[
\boxed{\lambda_{\max}=\frac{36}{5R_H}\simeq656\,\mathrm{nm}},
\qquad
\boxed{\lambda_{\min}=\frac4{R_H}\simeq364.6\,\mathrm{nm}}.
\]
严格数值应采用氢的\emph{有限核质量}Rydberg常数，
直接用 \(R_\infty\) 会有千分之一以下的小差别。

\paragraph{习题 28.2}
\(5\to1\) 的氢光子
\[
E_\gamma=13.6\,\mathrm{eV}
\left(1-\frac1{5^2}\right)\simeq13.06\,\mathrm{eV},
\quad\boxed{\lambda=\frac{hc}{E_\gamma}\simeq94.9\,\mathrm{nm}}.
\]
原子反冲动量与光子大小相同
\(p_H=h/\lambda\)，由 \(p_H=M_Hv\)：
\[
\boxed{v_{\rm recoil}=\frac{h}{M_H\lambda}
\simeq4.2\,\mathrm{m/s}}.
\]
更严格的能量守恒需从光子能量中减去
\(p_H^2/(2M_H)\)，但在上述有效数字下可忽略。

\paragraph{习题 28.3}
用Bohr半径
\(a_0=4\pi\epsilon_0\hbar^2/(m_ee^2)\)、
能级 \(E_n=-m_ee^4/(2(4\pi\epsilon_0)^2\hbar^2n^2)\)。
消去 \(e\)：
\[
\boxed{E_n=-\frac{\hbar^2}{2m_ea_0^2n^2}}.
\]
再用 \(\hbar^2/(m_ea_0)=e^2/(4\pi\epsilon_0)\)：
\[
\boxed{E_n=-\frac{e^2}{8\pi\epsilon_0a_0n^2}}.
\]
这些形式采用固定质子和经典Bohr的
非相对论库仑模型，不含精细结构。

\paragraph{习题 28.4}
\(n=1\) 时Bohr轨道半径 \(a_0\)，
由向心力条件
\(m_ev^2/a_0=e^2/(4\pi\epsilon_0a_0^2)\)。
因此动能 \(K=e^2/(8\pi\epsilon_0a_0)\)，
势能 \(U=-e^2/(4\pi\epsilon_0a_0)\)，
总能
\[
\boxed{E=K+U=-\frac{e^2}{8\pi\epsilon_0a_0}
=-\frac{\hbar^2}{2m_ea_0^2}=E_1}.
\]
相符是模型本身的内在一致性，不是对电子
具有确定经典轨道这一图像的证明。

\paragraph{习题 28.5}
由 \(hc/\lambda=E_{n_i}-E_{n_f}\)
和 \(E_n=-m_ee^4/[8\epsilon_0^2h^2n^2]\)，得
\[
\frac1\lambda
=\frac{m_ee^4}{8\epsilon_0^2h^3c}
\left(\frac1{n_f^2}-\frac1{n_i^2}\right).
\]
比较Balmer谱系 \(n_f=2\)：
\[
\boxed{R_\infty=\frac{m_ee^4}{8\epsilon_0^2h^3c}
\simeq1.0973732\times10^7\,\mathrm{m^{-1}}}.
\]
现代氢实际常数略小，原因是核的有限质量
须用约化质量 \(\mu\) 替代电子质量 \(m_e\)。

\paragraph{习题 28.6}
正电子和电子质量相同，约化质量
\(\mu=m_e/2\)，二者引力略去，
库仑Bohr半径以相对运动坐标定义为
\[
\boxed{r_{\min}=\frac{4\pi\epsilon_0\hbar^2}{\mu e^2}
=2a_0\simeq1.058\times10^{-10}\,\mathrm m}.
\]
其能级是氢的半数
\(E_n^{\rm Ps}\simeq-6.80/n^2\,\mathrm{eV}\)。
故 \(2\to1\) 发射
\(\Delta E=(3/4)6.80=5.10\,\mathrm{eV}\)，
\[
\boxed{\nu=\Delta E/h\simeq1.23\times10^{15}\,\mathrm{Hz}}.
\]
这里的轨道半径指两粒子\emph{相互距离}；
各粒子绕质心转动的半径是它的一半。

\paragraph{习题 28.7}
令太阳质量 \(M\)、地球质量 \(m\)，
圆轨道条件 \(mv^2/r=GMm/r^2\)；
角动量 \(L=mvr=m\sqrt{GMr}=n\hbar\)。
消去 \(v\)，得
\[
\boxed{r_n=\frac{n^2\hbar^2}{GMm^2}}.
\]
取 \(m=5.97\times10^{24}\,\mathrm{kg}\)、
\(M=1.99\times10^{30}\,\mathrm{kg}\)、
\(r=1.50\times10^{11}\,\mathrm m\)，
\[
\boxed{n\simeq\frac{m\sqrt{GMr}}{\hbar}
\simeq2.5\times10^{74}}.
\]
下一可能轨道的差
\[
\boxed{\Delta r=r_{n+1}-r_n
=r_n\left(\frac{2}{n}+\frac1{n^2}\right)
\simeq1.2\times10^{-63}\,\mathrm m}.
\]
这种宏观轨道的人为Bohr量子化给出
不可观测的间距，不意味着真实行星量子态
可以按原子轨道模型直接描述。

\paragraph{习题 28.8}
氢Balmer中两条最接近所测波长的谱线为
\(6\to2\) 的 \(H_\delta\)（约410.2nm）
和 \(5\to2\) 的 \(H_\gamma\)（约434.0nm）。
在低红移近似
\(z=(\lambda_{\rm obs}-\lambda_0)/\lambda_0
\simeq v/c\)。
于是
\[
\boxed{z_\delta\simeq\frac{411.7-410.2}{410.2}
\simeq3.66\times10^{-3}},\quad
\boxed{z_\gamma\simeq\frac{435.7-434.0}{434.0}
\simeq3.92\times10^{-3}}.
\]
结果因原书波长四位有效数字与谱线同位素/空气
波长表略不同而不完全一致，取平均
\[
\boxed{v\simeq1.1\times10^6\,\mathrm{m/s}
\simeq3.8\times10^{-3}c}.
\]
若用现代真空Rydberg值与观测空气波长
混合，会引入额外系统误差；
不宜冒充天文精密红移。

\paragraph{习题 28.9}
寿命 \(\tau=10^{-8}\,\mathrm s\) 的激发态
具有Lorentzian的自然线宽。
若用指数衰减偶极振幅
\(e^{-t/(2\tau)}\)，强度谱的频率\emph{全半高宽}
\[
\boxed{\Delta\nu_{\rm FWHM}=\frac1{2\pi\tau}
\simeq1.59\times10^7\,\mathrm{Hz}}.
\]
仅由粗略能量—时间不确定关系
\(\Delta E\,\Delta t\gtrsim\hbar/2\)
会得到典型半宽
\(\Delta\nu\sim1/(4\pi\Delta t)\simeq8.0\times10^6\,\mathrm{Hz}\)。
两者差因FWHM与半宽、寿命定义不同，
不得把它们无条件当成互相矛盾的测量。

\paragraph{习题 28.10}
取气体氢原子质量 \(m_H\)、
Maxwell一维速度分布，
\(\sigma_v=\sqrt{k_BT/m_H}\)。
低速多普勒偏移
\(\delta\nu\simeq-\nu_0v_z/c\)，
故频率标准差
\[
\boxed{\sigma_\nu
=\frac{\nu_0}{c}\sqrt{\frac{k_BT}{m_H}}
\simeq2.39\times10^9\,\mathrm{Hz}}
\]
(\(\nu_0=4.56\times10^{14}\,\mathrm{Hz},T=300\,\mathrm K\))。
高斯谱的FWHM
\[
\boxed{\Delta\nu_{\rm FWHM}
=2\sqrt{2\ln2}\,\sigma_\nu
\simeq5.63\times10^9\,\mathrm{Hz}}.
\]
若按源书S PDF376以\emph{最大热速度}
\(v_{\rm rms}=\sqrt{3k_BT/m_H}\)
取两侧频移，则全范围估计
\(2\nu_0v_{\rm rms}/c\simeq8.3\times10^9\,\mathrm{Hz}\)；
它与FWHM是\emph{不同的线宽定义}。

\paragraph{习题 28.11}
基态波函数
\(\psi_{100}=(\pi a_0^3)^{-1/2}e^{-r/a_0}\)。
径向概率密度
\[
P_r(r)=4\pi r^2|\psi_{100}|^2
=\frac{4r^2}{a_0^3}e^{-2r/a_0}.
\]
用 \(\int_0^\infty r^ne^{-\alpha r}dr=n!/\alpha^{n+1}\)：
\[
\boxed{\int_0^\infty P_r(r)\,dr
=\frac4{a_0^3}\frac{2!}{(2/a_0)^3}=1}.
\]
物理上意味着电子在整个空间被找到的总概率为1，
不是其位于特定Bohr轨道的概率为1。

\paragraph{习题 28.12}
\[
\langle r\rangle=\int_0^\infty rP_r(r)\,dr
=\frac4{a_0^3}\frac{3!}{(2/a_0)^4},
\qquad
\boxed{\langle r\rangle=\frac32a_0
\simeq7.94\times10^{-11}\,\mathrm m}.
\]
径向\emph{最概然}半径是 \(a_0\)，
与平均半径 \(1.5a_0\) 不同。

\paragraph{习题 28.13}
库仑势 \(U(r)=-e^2/(4\pi\epsilon_0r)\)：
\[
\langle U\rangle=-\frac{e^2}{4\pi\epsilon_0}
\int_0^\infty \frac1rP_r(r)\,dr
=-\frac{e^2}{4\pi\epsilon_0a_0},
\]
即
\[
\boxed{\langle U\rangle\simeq-27.2\,\mathrm{eV}}.
\]
总能 \(E_1=-13.6\,\mathrm{eV}\)，则
\[
\boxed{\langle K\rangle=E_1-\langle U\rangle
=+13.6\,\mathrm{eV}}.
\]
与经典圆轨道的总量一致，
但量子态的电子\emph{没有确定的经典轨迹}。

\paragraph{习题 28.14}
取源书三个规范化波函数的标准相位约定：
\[
\psi_{210}=\frac1{4\sqrt{2\pi a_0^3}}
\frac r{a_0}e^{-r/(2a_0)}\cos\theta,\quad
\psi_{21,\pm1}=
\frac1{8\sqrt{\pi a_0^3}}
\frac r{a_0}e^{-r/(2a_0)}
\sin\theta\,e^{\pm i\varphi}.
\]
令 \(Q(r)=(r/a_0)^2e^{-r/a_0}/a_0^3\)，
三个概率密度分别为
\[
\boxed{\rho_0=\frac{Q(r)\cos^2\theta}{32\pi}},
\qquad
\boxed{\rho_{+1}=\rho_{-1}
=\frac{Q(r)\sin^2\theta}{64\pi}}.
\]
因此 \(\rho_0+\rho_{+1}+\rho_{-1}
=Q(r)[\cos^2\theta+\sin^2\theta]/(32\pi)\)，
\[
\boxed{\rho_{\rm sum}=\frac{Q(r)}{32\pi}}
\]
\emph{与方向无关}。但独立的 \(m_l=0\)
态沿 \(z\) 轴最大、\(m_l=\pm1\) 态在赤道面最大，
并非各自球对称。
归一化：
\[
\int_{\mathbb R^3}|\psi_{210}|^2d^3r
=\frac1{32\pi a_0^5}
\underbrace{\int_0^\infty r^4e^{-r/a_0}dr}_{24a_0^5}
\underbrace{\int\cos^2\theta\,d\Omega}_{4\pi/3}
=\boxed{1}.
\]
上述是标准球谐函数的归一化；原书图28.5(b,c)
的三维绘制仍需逐图重绘核对。

\paragraph{习题 28.15}
轨道角动量大小
\(L=\hbar\sqrt{l(l+1)}=\sqrt2\hbar\)，
投影 \(L_z=m_l\hbar\)，
其中 \(m_l=-1,0,1\)。
因此
\[
\boxed{\cos\theta=\frac{m_l}{\sqrt2}},
\quad
\boxed{\theta=45^\circ,\ 90^\circ,\ 135^\circ}.
\]
注意这不是自旋角量子化；题设给 \(l=1\)。

\paragraph{习题 28.16}
该题指\emph{Lyman系最低频率双线}，
而非Balmer系。两条 \(2P\) 能级
裂距 \(\Delta E=4.5\times10^{-5}\,\mathrm{eV}\)，
相同 \(1S\) 末态，故
\[
\boxed{\Delta\nu=\frac{\Delta E}{h}
\simeq1.09\times10^{10}\,\mathrm{Hz}}.
\]
以Lyman-\(\alpha\) 波长
\(\lambda\simeq121.6\,\mathrm{nm}\) 计算小波长差：
\[
\boxed{|\Delta\lambda|
\simeq\frac{\lambda^2}{c}\Delta\nu
\simeq0.54\,\mathrm{pm}}.
\]
以源书简化等效自旋磁矩模型
\(\Delta E=2\mu_BB_{\rm eff}\)，
\[
\boxed{B_{\rm eff}=\Delta E/(2\mu_B)
\simeq0.39\,\mathrm T}.
\]
此“等效场”是模型对自旋轨道耦合的
简单参数化，不是原子核施加的均匀真实外场。

\paragraph{习题 28.17}
银原子外层有效自旋 \(s=\frac12\)，
\(l=0\)，\(|\boldsymbol S|=\sqrt{3/4}\hbar\)、
\(S_z=\pm\hbar/2\)，故
\[
\boxed{\cos\theta=\frac{S_z}{|\boldsymbol S|}
=\pm\frac1{\sqrt3}},\quad
\boxed{\theta\simeq54.74^\circ,\ 125.26^\circ}.
\]
电子磁矩与自旋反平行，
沿外场投影有 \(\mu_z\simeq\mp\mu_B\)，
故
\[
\boxed{U=-\mu_zB=\pm\mu_BB
\simeq\pm1.11\times10^{-23}\,\mathrm J
=\pm6.95\times10^{-5}\,\mathrm{eV}}.
\]
自旋角动量较接近外场方向的状态
磁势能较高；容易把磁矩和角动量方向混淆。

\paragraph{习题 28.18}
银原子的有效磁矩投影
\(\mu_z=\pm\mu_B\)。磁力
\[
\boxed{F_z=\mu_z\frac{dB}{dz}
=\pm(9.274\times10^{-24})(1.5\times10^3)
\simeq\pm1.39\times10^{-20}\,\mathrm N}.
\]
炉中速率平方取三维Maxwell平均
\(\bar v^2=3k_BT/m\)，在长 \(L=.040\,\mathrm m\)
梯度区内运动时间 \(L/\bar v\)；
两个分束在出口的间隔
\[
\boxed{\Delta z
=2\cdot\frac12\frac{\mu_B}{m}
\frac{dB}{dz}\left(\frac{L}{\bar v}\right)^2
=\frac{\mu_B L^2(dB/dz)}{3k_BT}
\simeq0.215\,\mathrm{mm}}.
\]
这忽略了磁场出口到记录板的附加漂移放大；
若有额外飞行距离，须另加横向位移，
不能照抄0.215mm当作任意屏距答案。

\paragraph{习题 28.19}
两核自旋取向能差
\(\Delta E=2\mu_zB\)：
\[
\boxed{\nu=\frac{2(1.41\times10^{-26})(1.60)}
{6.62607015\times10^{-34}}
\simeq6.81\times10^7\,\mathrm{Hz}=68.1\,\mathrm{MHz}},
\]
\[
\boxed{\lambda=c/\nu\simeq4.40\,\mathrm m}.
\]
共振频率属于射频范围。
核内局部场及化学位移近似忽略。

\paragraph{习题 28.20}
给定 \(n,l\)，允许
\(m_l=-l,-l+1,\ldots,l\) 共 \(2l+1\) 个，
每个轨道态自旋投影 \(m_s=\pm\frac12\)
可选两种；Pauli原理每种四量子数组合至多一电子，
\[
\boxed{N_{n,l}=2(2l+1)}.
\]
给定 \(n\)，有 \(l=0,\ldots,n-1\)，
\[
\boxed{N_n=2\sum_{l=0}^{n-1}(2l+1)=2n^2}.
\]
该结论是单电子原子轨道容量上限；
不等于多电子系统的具体能量次序。

\paragraph{习题 28.21}
基态电子构型如下：
\[
\begin{array}{ll}
\mathrm B(Z=5):&1s^2\,2s^2\,2p^1,\\
\mathrm{Ar}(18):&1s^2\,2s^2\,2p^6\,3s^2\,3p^6,\\
\mathrm{Cu}(29):&[\mathrm{Ar}]\,3d^{10}\,4s^1,\\
\mathrm{Br}(35):&[\mathrm{Ar}]\,3d^{10}\,4s^2\,4p^5.
\end{array}
\]
铜的 \(3d^{10}4s^1\) 是典型闭满d亚层例外，
不能机械写成 \(3d^94s^2\)。

\paragraph{习题 28.22}
加速电子动能 \(eU=30\,\mathrm{keV}\)，
Bremsstrahlung极短波端满足
\(eU=hc/\lambda_{\min}\)。
按题给波长 \(\lambda_{\min}=0.041\,\mathrm{nm}\)
反求
\[
\boxed{h=\frac{eU\lambda_{\min}}c
\simeq6.57\times10^{-34}\,\mathrm{J\,s}}.
\]
该结果与SI精确定义
\(h=6.62607015\times10^{-34}\,\mathrm{J\,s}\)
略有差别，因为题给截止波长只取两位有效数字。

\paragraph{习题 28.23}
短波极限 \(eU_{\min}=hc/\lambda\)，
\[
\boxed{U_{\min}
=\frac{hc}{e(0.100\times10^{-9}\,\mathrm m)}
\simeq1.24\times10^4\,\mathrm V=12.4\,\mathrm{kV}}.
\]
实际X射线管的辐射转换效率与加速电压下限
属于不同物理量。

\paragraph{习题 28.24}
四次碰撞电子依次损失
20keV、10keV、5keV、剩余5keV，
若全部分别发成一个光子，光子能量即这些数值。
使用 \(hc\simeq1.239842\,\mathrm{keV\,nm}\)，
\[
\boxed{\lambda_1\simeq0.0620\,\mathrm{nm}},\
\boxed{\lambda_2\simeq0.124\,\mathrm{nm}},\
\boxed{\lambda_3=\lambda_4\simeq0.248\,\mathrm{nm}}.
\]
实际制动辐射能谱是连续的，
题设是离散碰撞模型而非真实靶材的完整谱。

\paragraph{习题 28.25}
K\(\alpha\) 光子的能量：
\[
\boxed{\Delta E=\frac{hc}{3.16\,\mathrm{nm}}
\simeq392\,\mathrm{eV}}.
\]
按屏蔽近似的Moseley关系
\(\Delta E\simeq(3/4)\cdot13.6\,\mathrm{eV}
\,(Z-1)^2\)，所以
\[
Z\simeq1+\sqrt{\frac{392}{10.2}}
\simeq7.20,\quad
\boxed{Z\simeq7,\ \text{氮N}}.
\]
对轻元素而言，类氢屏蔽近似只用于粗判原子序数。

\paragraph{习题 28.26}
铜K\(\alpha\) 能量
\(\Delta E=(8.979-0.951)\,\mathrm{keV}
=8.028\,\mathrm{keV}\)，
\[
\lambda=\frac{hc}{\Delta E}\simeq0.15444\,\mathrm{nm}.
\]
原题给的是\emph{相对晶面}的掠射角
\(\varphi=74.1^\circ\)，一级Bragg条件
\(2d\sin\varphi=\lambda\)：
\[
\boxed{d=\frac{\lambda}{2\sin74.1^\circ}
\simeq0.0803\,\mathrm{nm}}.
\]
这是题设指定的\emph{Na离子平面间距}的推算值；
不要未经查验晶面指数便把这一结果
等同于NaCl常见的宏观晶胞常数。

\paragraph{习题 28.27}
激光光子能量
\[
\boxed{\Delta E=\frac{hc}{10.6\,\mathrm{\mu m}}
\simeq0.117\,\mathrm{eV}}.
\]
对非简并且热平衡的两能级，Boltzmann比值
\[
\boxed{\frac{N_{\rm up}}{N_{\rm low}}
=\exp(-\Delta E/(k_BT))
\simeq0.01084\ (300\,\mathrm K)}
\]
即上能级约为下能级的1.08\%。
若反转为 \(N_{\rm up}/N_{\rm low}=1.01\)，
形式温度满足
\[
\boxed{T_{\rm eff}
=-\frac{\Delta E}{k_B\ln1.01}
\simeq-1.36\times10^5\,\mathrm K}.
\]
负温度只适用于有上限、并处于
反转分布的\emph{特定两能级子系统}，
不是全部气体分子的平动温度变负。

\paragraph{习题 28.28}
脉冲持续 \(\Delta t=10\,\mathrm{fs}=10^{-14}\,\mathrm s\)。
光学周期 \(T=\lambda/c\)；因此
\[
\boxed{N=\frac{\Delta t}{T}
=\frac{c\Delta t}{\lambda}
=\frac{(3.00\times10^8)(10^{-14})}{500\times10^{-9}}
\simeq6}.
\]
极短脉冲只约6个光学周期，不能使用严格
单色无限时长的平面波模型描述带宽。

\paragraph{习题 28.29}
脉冲持续 \(\tau=12\,\mathrm{ps}\)，
\(\lambda=694.4\,\mathrm{nm}\)，能量0.150J。
空间长度：
\[
\boxed{\ell=c\tau
\simeq3.60\times10^{-3}\,\mathrm m
=3.60\,\mathrm{mm}}.
\]
平均脉冲功率
\[
\boxed{P_{\rm pulse}=\frac{0.150}{12\times10^{-12}}
=1.25\times10^{10}\,\mathrm W=12.5\,\mathrm{GW}}.
\]
光子数
\[
\boxed{N_\gamma=\frac{W}{hc/\lambda}
=\frac{W\lambda}{hc}\simeq5.24\times10^{17}}.
\]
\textbf{配套S PDF382疑似单位换算笔误：}
它写 \(3.6\times10^{-3}\,\mathrm m=0.36\,\mathrm{mm}\)，
但正确换算是\emph{3.6mm}，差10倍。
源书数值不应未经复核直接纳入答案库。

\paragraph{习题 28.30}
体积 \(200\,\mathrm{\mu m^3}\) 不参与
光子发射\emph{总速率}计算。
每个光子能量 \(hc/\lambda\)，
连续光功率 \(P=5.0\,\mathrm{mW}\)：
\[
\boxed{\dot N_\gamma=\frac{P\lambda}{hc}
=\frac{(.005)(.80\times10^{-6})}
{(6.62607015\times10^{-34})(2.99792458\times10^8)}
\simeq2.01\times10^{16}\,\mathrm{s^{-1}}}.
\]
若问单位体积光子密度，才需补充模式寿命或
腔内循环功率等信息，不能仅据200μm$^3$推算。

\paragraph{习题 28.31}
入射圆光束半径 \(r_0=D/2=1.50\,\mathrm{mm}\)：
\[
\boxed{I_0=\frac{P}{\pi r_0^2}
=\frac5{\pi(.0015)^2}
\simeq7.07\times10^5\,\mathrm{W/m^2}}.
\]
圆孔Airy第一零点角度
\(\theta_1\simeq1.22\lambda/D\)，
焦平面中心斑半径
\[
\boxed{R\simeq\frac{1.22\lambda f}{D}
=\frac{1.22(515\times10^{-9})(.035)}{.003}
\simeq7.33\,\mathrm{\mu m}}.
\]
假定斑内占总能量84\%，
中心斑的\emph{区域平均强度}
\[
\boxed{\bar I_{\rm Airy}
=\frac{.84P}{\pi R^2}
\simeq2.49\times10^{10}\,\mathrm{W/m^2}}.
\]
Airy亮斑中央的\emph{局部峰值强度}不等于这里的
圆域平均值，题中“强度”按配套书区域平均解释。

\paragraph{习题 28.32}
刚性对称双原子转子的转动谱线
\(E_J=\hbar^2J(J+1)/(2I)\)，
相邻\(J\to J-1\)谱线的相邻间隔
\(\Delta\nu=\hbar/(\pi I)\)。
两个相同质量 \(m\) 原子距离 \(r\)，
对质心转动的惯量 \(I=2m(r/2)^2=mr^2/2\)。
于是
\[
\boxed{r=\sqrt{\frac{\hbar}{\pi m\Delta\nu}}
=\sqrt{\frac{1.05457\times10^{-34}}
{\pi(2.66\times10^{-26})(8.6\times10^{10})}}
\simeq1.21\times10^{-10}\,\mathrm m
=0.121\,\mathrm{nm}}.
\]
若把单个原子质量 \(m\) 误当总体系约化质量，
惯量会差2倍。

\paragraph{习题 28.33}
此题不是把电子绕质子的单质点轨道当作分子转子，
而是\emph{氢分子绕两原子中心连线的自身轴转动}。
质子在自旋轴附近，其转动惯量可忽略；
每个氢原子的电子云近似均匀实心球，
绕其直径的惯量为 \(I_{\rm atom}=(2/5)m_ea_0^2\)。
两个球的惯量为
\[
I=2I_{\rm atom}=\frac45m_ea_0^2.
\]
刚性转子第一激发态 \(J=1\) 相对 \(J=0\)：
\[
\boxed{\Delta E=E_1-E_0=\frac{\hbar^2}{I}
=\frac{5\hbar^2}{4m_ea_0^2}
\simeq34\,\mathrm{eV}}.
\]
等价温度 \(\Delta E/k_B\simeq3.9\times10^5\,\mathrm K\)，
比室温热能 \(k_BT\simeq0.026\,\mathrm{eV}\) 大许多，
\[
\boxed{\text{该沿分子轴旋转模式在常温冻结，
对氢气比热无显著贡献}}.
\]
这不同于真实H$_2$绕垂直于分子轴的低能量转动谱；
两种转轴不可混淆。

\paragraph{习题 28.34}
CO的两原子约化质量为
\(\mu=m_Cm_O/(m_C+m_O)
=(12\cdot16/28)u\simeq1.139\times10^{-26}\,\mathrm{kg}\)。
谐振动频率
\(\nu=(2\pi)^{-1}\sqrt{k/\mu}\)，故
\[
\boxed{k=(2\pi\nu)^2\mu
=[2\pi(6.42\times10^{13})]^2
(1.139\times10^{-26})
\simeq1.85\times10^3\,\mathrm{N/m}}.
\]
此为分子势能在平衡键长附近的二阶曲率，
不是分子间宏观弹簧常数。
